\subsection{自同态的特征多项式}

设 M 是维数为 n 的自由 A-模，u 是 M 的一个自同态。考虑两个不定元的多项式环 A{[}X, Y{]} 以及 A{[}X, Y{]}-模 M{[}X, Y{]} = A{[}X, Y{]} ⊗\_A M；令 \(\bar{u}\) 为由 u 典范地导出的 A{[}X, Y{]}-模 M{[}X, Y{]} 的自同态（II，§5，第1号）。由第 5 号命题 11 可知，

\[
\det (\mathbf {X} - \mathbf {Y} \bar {\boldsymbol {u}}) = \sum_ {j = 0} ^ {n} (- 1) ^ {j} \operatorname{Tr} \left(\bigwedge^ {j} (u)\right) \mathbf {X} ^ {n - j} \mathbf {Y} ^ {j}\tag{49}
\]

这是因为，若 \(U\) 是 \(u\) 关于某个基 \((e_i)_{1 \leqslant i \leqslant n}\) （ \(\mathbf{M}\) 的）的矩阵，则 \(U\) 就是 \(\bar{u}\) 关于基 \((1 \otimes e_i)_{1 \leqslant i \leqslant n}\) （ \(\mathbf{M}[\mathbf{X}, \mathbf{Y}]\) 的）的矩阵，因此

\[
\operatorname{Tr} \left(\bigwedge^ {j} (\bar {u})\right) = \operatorname{Tr} \left(\bigwedge^ {j} (u)\right).
\]

定义 3. 设 M 为有限维自由 A-模，u 为 M 的一个自同态。自同态 X - \(\bar{u}\) （自由 A{[}X{]}-模 M{[}X{]} 的）的行列式称为 u 的特征多项式，记作 \(\chi_u(X)\) 。

若 \(\mathbf{M}\) 的秩为 \(n\) ，则由 (49) 可知

\[
\chi_ {u} (\mathbf {X}) = \sum_ {j = 0} ^ {n} (- 1) ^ {j} \operatorname{Tr} \left(\bigwedge^ {j} (u)\right) \mathbf {X} ^ {n - j}\tag{50}
\]

因为 \(\operatorname{det}(\mathbf{X} - \mathbf{Y}\bar{u}) = \operatorname{det}(\mathbf{X}.I_n + \mathbf{Y}U)\) 且 \(\operatorname{det}(\mathbf{X} - \bar{u}) = \operatorname{det}(\mathbf{X}.I_n - U)\) 。由此可以看出 \(\chi_u(\mathbf{X})\) 是一个次数为 \(n\) 的首一多项式，其中 \(\mathbf{X}^{n-1}\) 的系数是 \(-\mathrm{Tr}(u)\) ，常数项是 \((-1)^n \operatorname{det}(u)\) 。

命题 20（”凯莱-哈密顿定理”）。对于有限维自由 A-模的每个自同态 \(u\) ，有 \(\chi_u(u) = 0\) 。

在命题 18（§ 3，第 10 号）的记号下，对于所有 \(x \in \mathbf{M}\) ， \(\chi_u(u)(x)\) 是在映射 \(\phi\) 下 \(\chi_u(\mathbf{X}) \otimes x\) 的像。但若 \(v\) 是 \(\mathbf{M}[\mathbf{X}]\) 的自同态、即 \(\mathbf{X} - \bar{u}\) 的余转置（第 6 号），那么

\[
\chi_ {u} (\mathbf {X}) \otimes x = \chi_ {u} (\mathbf {X}) (1 \otimes x) = (\mathbf {X} - \bar {u}) (v (1 \otimes x))
\]

且结论由第 10 号的命题 18 得出。

